\par
\SourceLeaf{391}%
Let \(r\). be the distance of the centre of the desired ellipse from the origin of the abscissas, and \(\varphi\) the angle that the straight line drawn from the latter point to the former makes with the axis of abscissas. In other words, let \(r\cos\varphi,r\sin\varphi\) be the coordinates of the centre of the ellipse. From this it is easy to find that the perpendicular from this centre to the first side of the quadrilateral has length
\[
=a-r\cos(A-\varphi)
\]
The perpendiculars to the other three sides are expressed in the same way.

If the semimajor axis of the ellipse is denoted by \(\alpha\), the semiminor axis by \(\beta\), and the inclination of the latter to the axis of abscissas by \(\psi\), then clearly \(A-\psi\) is the inclination to the semiminor axis of the perpendicular from the centre to the first side of the quadrilateral. If that side is to be tangent to the ellipse, this perpendicular is expressed, for familiar reasons, by
\[
\BareRoot{[\alpha\alpha\sin(A-\psi)^2+\beta\beta\cos(A-\psi)^2]}
\]
We thus have the equation
\[
a-r\cos(A-\varphi)=\BareRoot{[\alpha\alpha\sin(A-\psi)^2+\beta\beta\cos(A-\psi)^2]}
\]
and three more of the same form, obtained by replacing \(a\) and \(A\) with the symbols for the other sides. If we therefore remove the radical and set, for brevity,
\[
\begin{aligned}
rr-\alpha\alpha-\beta\beta&=t\\
\alpha\alpha-\beta\beta&=u
\end{aligned}
\]
then our four equations are
\begingroup\fontsize{10}{12}\selectfont
\[
\begin{array}{rl}
\mathrm{I.}&2aa+t-4ar\cos(A-\varphi)+rr\cos2(A-\varphi)-u\cos2(A-\psi)=0\\
\mathrm{II.}&2a'a'+t-4a'r\cos(A'-\varphi)+rr\cos2(A'-\varphi)-u\cos2(A'-\psi)=0\\
\mathrm{III.}&2a''a''+t-4a''r\cos(A''-\varphi)+rr\cos2(A''-\varphi)-u\cos2(A''-\psi)=0\\
\mathrm{IV.}&2a'''a'''+t-4a'''r\cos(A'''-\varphi)+rr\cos2(A'''-\varphi)-u\cos2(A'''-\psi)=0
\end{array}
\]
\endgroup
If we multiply the first equation by \(\sin2(A''-A')\), the second by \(\sin2(A-A'')\), and the third by \(\sin2(A'-A)\), and add the products, we obtain (see Art.\,78 of my \textit{Theoria motus corporum coelestium})
